% Written by `lab paper tactic-found`; do not edit by hand.
\begin{tabularx}{\textwidth}{lL}
\toprule
Mark & The squares it is set on \\
\midrule
\texttt{en\_passant} & The square a pawn may be taken on in passing. \\
\texttt{setup\_from} & The square the opponent's move left: the move that set the position up. \\
\texttt{setup\_to} & The square that move reached. \\
\texttt{solver1\_from} & The square the line's first move leaves. \\
\texttt{solver1\_to} & The square the line's first move reaches. \\
\texttt{reply1\_from} & The square the first reply leaves. \\
\texttt{reply1\_to} & The square the first reply reaches. \\
\texttt{solver2\_from} & The square the solver's second move leaves. \\
\texttt{solver2\_to} & The square the solver's second move reaches. \\
\texttt{reply2\_from} & The square the second reply leaves. \\
\texttt{reply2\_to} & The square the second reply reaches. \\
\texttt{solver3\_from} & The square the solver's third move leaves. \\
\texttt{solver3\_to} & The square the solver's third move reaches. \\
\texttt{later} & Every square a later move of the line leaves or reaches. \\
\bottomrule
\end{tabularx}
